% GENERATED FILE. Edit canonical lessons, then run npm run generate:books.
% canonical-source-hash: fnv1a64:5c94ca6a6cb439ac
% canonical-lessons: KA-C421-bedike, KA-C421-puraisu, KA-C421-korate, KA-C421-raphtu, KA-C421-hanadubbara

\chapter{Supply and Demand}
\label{ch:ka-supply-and-demand}
\hlchaptermodality{421}

\begin{chapteropening}
\textbf{By the end of this chapter:} \emph{I can say a demand, to supply, a shortage, export and inflation.}

Complete the last lesson of chapter 421: I can say a demand, to supply, a shortage, export and inflation.
\end{chapteropening}

\section[\texorpdfstring{\kn{ಬೇಡಿಕೆ} (bēḍike)}{ಬೇಡಿಕೆ (bēḍike)}]{\kn{ಬೇಡಿಕೆ} (bēḍike) — a demand}

\label{lesson:KA-C421-bedike}



\begin{quote}
Before the new one: say the Kannada for uncertain, then the Kannada for a calamity.
\end{quote}

\subsection*{You'll want to know: \kn{ಬೇಡಿಕೆ}}
\textbf{\kn{ಬೇಡಿಕೆ}} — \emph{bēḍike} — \textquotedblleft{}a demand\textquotedblright{}.

\begin{quote}
\kn{ಹಬ್ಬದ} \kn{ಸಮಯದಲ್ಲಿ} \kn{ಹೂವಿಗೆ} \kn{ಬೇಡಿಕೆ} \kn{ಹೆಚ್ಚು}.
\end{quote}

\begin{quote}
\emph{habbada samayadalli hūvige bēḍike heccu.}
\end{quote}

\begin{quote}
At festival time flowers are in high demand.
\end{quote}

Also a demand made by people: \kn{ಕಾರ್ಮಿಕರ} \kn{ಬೇಡಿಕೆ}.

\begin{cousinweb}
\textbf{\kn{ಬೇಡಿಕೆ}} (\emph{bēḍike}) is native Dravidian, a noun from \emph{bēḍu}, \textquotedblleft{}to ask for\textquotedblright{}, listed with \emph{bēḍike} in DEDR 5528.
\end{cousinweb}

\begin{grammarlens}[title={What you've built}]
One more word for saying what happens if.
\end{grammarlens}

\subsection*{Your turn}
\begin{itemize}
\raggedright
  \item \emph{Say it:} \emph{bēḍike}
  \item \emph{Say it:} \emph{bēḍike}, once more
  \item \emph{Say it:} \emph{habbada samayadalli hūvige bēḍike heccu.} — At festival time flowers are in high demand.
  \item \emph{Recall:} say the Kannada for uncertain, then the Kannada for a calamity, then say \emph{bēḍike} again
\end{itemize}

\begin{tcolorbox}[breakable,colback=teal!4,colframe=teal!35!black,title={Before you move on}]
What does \kn{ಬೇಡಿಕೆ} mean? (\textquotedblleft{}A demand\textquotedblright{}.) Say it once more.
\end{tcolorbox}
\section[\texorpdfstring{\kn{ಪೂರೈಸು} (pūraisu)}{ಪೂರೈಸು (pūraisu)}]{\kn{ಪೂರೈಸು} (pūraisu) — to supply}

\label{lesson:KA-C421-puraisu}



\begin{quote}
Before the new one: say the Kannada for a calamity, then the Kannada for a demand.
\end{quote}

\subsection*{You'll want to know: \kn{ಪೂರೈಸು}}
\textbf{\kn{ಪೂರೈಸು}} — \emph{pūraisu} — \textquotedblleft{}to supply\textquotedblright{}.

\begin{quote}
\kn{ಆ} \kn{ಹಳ್ಳಿ} \kn{ನಗರಕ್ಕೆ} \kn{ಹಾಲು} \kn{ಪೂರೈಸುತ್ತದೆ}.
\end{quote}

\begin{quote}
\emph{ā haḷḷi nagarakke hālu pūraisuttade.}
\end{quote}

\begin{quote}
That village supplies milk to the city.
\end{quote}

Past (I …): \kn{ಪೂರೈಸಿದೆ}, \emph{pūraiside}.

Also to complete: \kn{ಓದು} \kn{ಪೂರೈಸು}, to finish one's studies.

\begin{cousinweb}
\textbf{\kn{ಪೂರೈಸು}} (\emph{pūraisu}) is built on Sanskrit \emph{pūra}, \textquotedblleft{}filling\textquotedblright{} (CDIAL 8330), with a Kannada verb ending: to make full.
\end{cousinweb}

\begin{grammarlens}[title={What you've built}]
One more word for saying what happens if.
\end{grammarlens}

\subsection*{Your turn}
\begin{itemize}
\raggedright
  \item \emph{Say it:} \emph{pūraisu}
  \item \emph{Say it:} \emph{pūraisu}, once more
  \item \emph{Say it:} \emph{ā haḷḷi nagarakke hālu pūraisuttade.} — That village supplies milk to the city.
  \item \emph{Recall:} say the Kannada for a calamity, then the Kannada for a demand, then say \emph{pūraisu} again
\end{itemize}

\begin{tcolorbox}[breakable,colback=teal!4,colframe=teal!35!black,title={Before you move on}]
What does \kn{ಪೂರೈಸು} mean? (\textquotedblleft{}To supply\textquotedblright{}.) Say it once more.
\end{tcolorbox}
\section[\texorpdfstring{\kn{ಕೊರತೆ} (korate)}{ಕೊರತೆ (korate)}]{\kn{ಕೊರತೆ} (korate) — a shortage; a lack}

\label{lesson:KA-C421-korate}



\begin{quote}
Before the new one: say the Kannada for a demand, then the Kannada for to supply.
\end{quote}

\subsection*{You'll want to know: \kn{ಕೊರತೆ}}
\textbf{\kn{ಕೊರತೆ}} — \emph{korate} — \textquotedblleft{}a shortage; a lack\textquotedblright{}.

\begin{quote}
\kn{ಬೇಸಿಗೆಯಲ್ಲಿ} \kn{ನೀರಿನ} \kn{ಕೊರತೆ} \kn{ಇರುತ್ತದೆ}.
\end{quote}

\begin{quote}
\emph{bēsigeyalli nīrina korate iruttade.}
\end{quote}

\begin{quote}
In summer there is a shortage of water.
\end{quote}

What is lacking takes -\kn{ದ} / -\kn{ಇನ}: \kn{ಹಣದ} \kn{ಕೊರತೆ}.

\begin{cousinweb}
\textbf{\kn{ಕೊರತೆ}} (\emph{korate}) is native Dravidian: DEDR 1851, the \textquotedblleft{}short, small\textquotedblright{} family, lists Kannada \emph{koṟate}, \textquotedblleft{}a deficiency\textquotedblright{}.
\end{cousinweb}

\begin{grammarlens}[title={What you've built}]
One more word for saying what happens if.
\end{grammarlens}

\subsection*{Your turn}
\begin{itemize}
\raggedright
  \item \emph{Say it:} \emph{korate}
  \item \emph{Say it:} \emph{korate}, once more
  \item \emph{Say it:} \emph{bēsigeyalli nīrina korate iruttade.} — In summer there is a shortage of water.
  \item \emph{Recall:} say the Kannada for a demand, then the Kannada for to supply, then say \emph{korate} again
\end{itemize}

\begin{tcolorbox}[breakable,colback=teal!4,colframe=teal!35!black,title={Before you move on}]
What does \kn{ಕೊರತೆ} mean? (\textquotedblleft{}A shortage; a lack\textquotedblright{}.) Say it once more.
\end{tcolorbox}
\section[\texorpdfstring{\kn{ರಫ್ತು} (raphtu)}{ರಫ್ತು (raphtu)}]{\kn{ರಫ್ತು} (raphtu) — export}

\label{lesson:KA-C421-raphtu}



\begin{quote}
Before the new one: say the Kannada for to supply, then the Kannada for a shortage.
\end{quote}

\subsection*{You'll want to know: \kn{ರಫ್ತು}}
\textbf{\kn{ರಫ್ತು}} — \emph{raphtu} — \textquotedblleft{}export\textquotedblright{}.

\begin{quote}
\kn{ನಮ್ಮ} \kn{ಕಾಫಿ} \kn{ಬೇರೆ} \kn{ದೇಶಗಳಿಗೆ} \kn{ರಫ್ತು} \kn{ಆಗುತ್ತದೆ}.
\end{quote}

\begin{quote}
\emph{namma kāphi bēre dēśagaḷige raphtu āguttade.}
\end{quote}

\begin{quote}
Our coffee is exported to other countries.
\end{quote}

\kn{ರಫ್ತು} \kn{ಮಾಡು}, to export. \kn{ಆಮದು}, import.

\begin{cousinweb}
\textbf{\kn{ರಫ್ತು}} (\emph{raphtu}) is usually traced to Persian \emph{raft}, \textquotedblleft{}gone\textquotedblright{}, through Hindustani \emph{raftanī}, \textquotedblleft{}export\textquotedblright{}.
\end{cousinweb}

\begin{grammarlens}[title={What you've built}]
One more word for saying what happens if.
\end{grammarlens}

\subsection*{Your turn}
\begin{itemize}
\raggedright
  \item \emph{Say it:} \emph{raphtu}
  \item \emph{Say it:} \emph{raphtu}, once more
  \item \emph{Say it:} \emph{namma kāphi bēre dēśagaḷige raphtu āguttade.} — Our coffee is exported to other countries.
  \item \emph{Recall:} say the Kannada for to supply, then the Kannada for a shortage, then say \emph{raphtu} again
\end{itemize}

\begin{tcolorbox}[breakable,colback=teal!4,colframe=teal!35!black,title={Before you move on}]
What does \kn{ರಫ್ತು} mean? (\textquotedblleft{}Export\textquotedblright{}.) Say it once more.
\end{tcolorbox}
\section[\texorpdfstring{\kn{ಹಣದುಬ್ಬರ} (haṇadubbara)}{ಹಣದುಬ್ಬರ (haṇadubbara)}]{\kn{ಹಣದುಬ್ಬರ} (haṇadubbara) — inflation}

\label{lesson:KA-C421-hanadubbara}



\begin{quote}
Before the new one: say the Kannada for a shortage, then the Kannada for export.
\end{quote}

\subsection*{You'll want to know: \kn{ಹಣದುಬ್ಬರ}}
\textbf{\kn{ಹಣದುಬ್ಬರ}} — \emph{haṇadubbara} — \textquotedblleft{}inflation\textquotedblright{}.

\begin{quote}
\kn{ಹಣದುಬ್ಬರದಿಂದ} \kn{ಎಲ್ಲದರ} \kn{ಬೆಲೆ} \kn{ಏರಿತು}.
\end{quote}

\begin{quote}
\emph{haṇadubbaradinda elladara bele ēritu.}
\end{quote}

\begin{quote}
With inflation the price of everything went up.
\end{quote}

Word by word: money swelling.

\begin{cousinweb}
\textbf{\kn{ಹಣದುಬ್ಬರ}} (\emph{haṇadubbara}) is a modern Kannada word built from \emph{haṇa}, \textquotedblleft{}money\textquotedblright{}, and \emph{ubbara}, \textquotedblleft{}a swelling\textquotedblright{}, from native \emph{ubbu}, \textquotedblleft{}to swell\textquotedblright{} (DEDR 666).
\end{cousinweb}

\begin{grammarlens}[title={What you've built}]
One more word for saying what happens if.
\end{grammarlens}

\subsection*{Your turn}
\begin{itemize}
\raggedright
  \item \emph{Say it:} \emph{haṇadubbara}
  \item \emph{Say it:} \emph{haṇadubbara}, once more
  \item \emph{Say it:} \emph{haṇadubbaradinda elladara bele ēritu.} — With inflation the price of everything went up.
  \item \emph{Recall:} say the Kannada for a shortage, then the Kannada for export, then say \emph{haṇadubbara} again
\end{itemize}

\begin{tcolorbox}[breakable,colback=teal!4,colframe=teal!35!black,title={Before you move on}]
What does \kn{ಹಣದುಬ್ಬರ} mean? (\textquotedblleft{}Inflation\textquotedblright{}.) Say it once more.
\end{tcolorbox}
